\section{Excitation of Free Carriers}

Current experimental methods for generating femtosecond lasers in the near ultraviolet range (NUV) remain complex, which limits research on the linear optical transitions of wide bandgap dielectric materials such as $\mathrm{SiO}_2$.

Given these experimental limitations, research has focused mainly on the study of nonlinear absorption processes that generate free carriers in the conduction band. By responding to the electric field, these electrons increase their energy and trigger effects such as avalanche ionization, which modify the optical properties of the material, in particular its absorbance~\cite{mao2004, bulgakova2010}.

In this section, we explain the nonlinear mechanisms responsible for the generation of free carriers, that is, the excitation of electrons from the valence band to the conduction band. In wide bandgap dielectrics, these carriers are central to the light matter interaction, since they are the only way to transfer optical energy to the crystal lattice of the material, which is otherwise transparent to low intensity light.
